% =====================================================================
%  Lecture 10 -- Diameter Protocol                        (English edition)
% =====================================================================
\documentclass[11pt,a4paper]{article}
\providecommand\ACNroot{../..}
\input{\ACNroot/templates/lecture-en}

\ACNset{number}{10}
\ACNset{title}{Diameter Protocol}
\ACNset{shorttitle}{Diameter}
\ACNset{titleother}{پروتکل Diameter}
\ACNset{part}{Part II -- 4G access and core (Lectures 5--13)}
\ACNset{time}{6--8 hours (3 h reading, 3--5 h exercises)}
\ACNset{level}{Intermediate}
\ACNset{prereq}{Lectures 5 and 9; TCP/SCTP; binary protocol encoding (TLV)}
\ACNset{gen}{4}

\begin{document}
\ACNcover

% ---------------------------------------------------------------------
\section{Learning objectives}
After studying this lecture you should be able to:
\begin{objectives}
  \item \textbf{Explain} the limitations of RADIUS that Diameter was designed to remove.
  \item \textbf{Describe} Diameter nodes and agents, peer connections and realm-based routing.
  \item \textbf{Decode} a Diameter header and AVP, including flags, vendor IDs and padding, and
        \textbf{compute} message lengths.
  \item \textbf{List} the base commands and result-code classes, and \textbf{trace} the S6a, Gx,
        Rx and Gy applications used in LTE.
  \item \textbf{Size} a credit-control quota and \textbf{estimate} the transaction load of an HSS.
  \item \textbf{Explain} why 5G replaced Diameter with HTTP/2 service-based interfaces, and how the two
        interwork.
\end{objectives}

% ---------------------------------------------------------------------
\section{Prerequisite knowledge}
You need the EPC reference points of Lectures~5 and~9 (S6a, Gx, Rx, Gy) and a working knowledge of
transport protocols: TCP's reliable byte stream and SCTP's message-oriented, multi-homed associations
(Lecture~6). Binary protocols are usually built from \emph{type--length--value} (TLV) elements; Diameter's
AVP is one.

\paragraph{Key terms} The vocabulary of this lecture with its Persian equivalents:
\begin{keyterms}
AAA & \fa{احراز هویت، مجوزدهی، حساب‌داری} & Authentication, authorisation and accounting framework. \\
AVP & \fa{جفت ویژگی--مقدار} & Diameter's TLV information element. \\
Command code & \fa{کد فرمان} & Identifies a request/answer pair (e.g.\ 316 = Update-Location). \\
Application-ID & \fa{شناسهٔ کاربرد} & Identifies the interface (e.g.\ 16777251 = S6a/S6d). \\
Realm & \fa{قلمرو} & Administrative domain used for routing (e.g.\ an operator's domain). \\
Relay / proxy agent & \fa{عامل رله / پراکسی} & Agents that forward messages (proxy may modify them). \\
DRA / DEA & \fa{عامل مسیریابی / لبهٔ Diameter} & Routing agent inside, edge agent at the border of a network. \\
CER / DWR / DPR & \fa{تبادل قابلیت، پاسبان، قطع} & Base peer-management commands. \\
Result-Code & \fa{کد نتیجه} & Outcome of a request (2001 = success). \\
Credit control & \fa{کنترل اعتبار} & Online charging with quota grants (Gy). \\
\end{keyterms}

% ---------------------------------------------------------------------
\section{Introduction and motivation}
Every time an LTE phone attaches, an MME asks the HSS for authentication vectors and the subscriber's
profile; every time a VoLTE call starts, the IMS asks the PCRF for a guaranteed bearer; every few
megabytes, a prepaid subscriber's P-GW asks the charging system for more credit. All these dialogues
-- between the control nodes of the core, never touching user data -- use one protocol: Diameter
\cite{rfc6733}.

Diameter descends from RADIUS \cite{rfc2865}, the dial-up era's AAA protocol (the name is a pun: a
diameter is twice a radius). RADIUS could not meet operator-grade requirements -- reliability,
large attributes, server-initiated messages, roaming between realms -- so the IETF designed Diameter as an
extensible base protocol plus \emph{applications}. 3GPP adopted it for IMS (Release~5) and made it the
signalling backbone of the EPC. In 5G it is being replaced by HTTP/2 APIs, but Diameter networks will
carry 4G traffic for many years, and the design concepts -- realms, agents, capability exchange -- recur
in the 5G service-based architecture.

% ---------------------------------------------------------------------
\section{Fundamental concepts}

\subsection{From RADIUS to Diameter}
\begin{table}[htbp]
\centering\small
\caption{RADIUS versus Diameter.}\label{tab:radius}
\rowcolors{2}{acnmist}{white}
\begin{tabularx}{\linewidth}{@{}l Y Y@{}}
\toprule
\rowcolor{white} & RADIUS & Diameter \\
\midrule
Transport              & UDP (application-level retransmission) & TCP or SCTP (reliable), port 3868 \\
Attribute length       & 8-bit length field, at most 253 bytes of data & 24-bit length field \\
Outstanding requests   & 8-bit identifier: 256 per client--server pair & 32-bit hop-by-hop and end-to-end identifiers \\
Server-initiated messages & Not supported (only extensions)     & Supported (e.g.\ Cancel-Location, Re-Auth) \\
Peer liveness / failover & No standard mechanism              & Watchdog (DWR/DWA) and failover to alternative peers \\
Capability negotiation & None                                   & CER/CEA exchange supported applications \\
Agents                 & Proxies, ad hoc                        & Relay, proxy, redirect and translation agents \\
Routing                & Static, by server address              & By Destination-Realm and Destination-Host \\
Security               & Shared secret, MD5-based hiding         & TLS/DTLS or IPsec \\
\bottomrule
\end{tabularx}
\end{table}

\subsection{Nodes and agents}
A Diameter \emph{client} (e.g.\ the MME) sends requests to a \emph{server} (e.g.\ the HSS). Between
them, \emph{agents} may forward messages:
\begin{itemize}
  \item a \textbf{relay} forwards based on routing information without inspecting the application;
  \item a \textbf{proxy} forwards and may modify messages to enforce policy;
  \item a \textbf{redirect} agent answers with the address of the right server instead of forwarding;
  \item a \textbf{translation} agent converts to another protocol (e.g.\ RADIUS).
\end{itemize}
Operators deploy \textbf{Diameter Routing Agents} (DRAs) as central signalling routers, so that each
client needs connections to two DRAs instead of a full mesh, and \textbf{Diameter Edge Agents} (DEAs) at
the border towards roaming partners (via an IP exchange, IPX), hiding the internal topology.

\subsection{Peers and realms}
Two directly connected nodes are \emph{peers}. Before exchanging application messages they perform a
\emph{capabilities exchange} (CER/CEA) announcing their identity (Origin-Host, Origin-Realm) and supported
applications. Each node keeps a \emph{peer table} and a \emph{realm routing table}. A request carries
Destination-Realm (mandatory for routed requests) and optionally Destination-Host; nodes forward it to a
peer that serves that realm and application. In 3GPP networks realms follow the pattern
\texttt{epc.mnc<MNC>.mcc<MCC>.3gppnetwork.org} \cite{ts23003}, so a visited MME can route an S6a request to
a subscriber's home HSS using only the IMSI.

% ---------------------------------------------------------------------
\section{Detailed technical explanation}

\subsection{Message structure}
A Diameter message is a 20-byte header followed by AVPs (\cref{fig:header}) \cite{rfc6733}.
\begin{itemize}
  \item \textbf{Version} = 1; \textbf{Message Length} includes the header and is a multiple of 4.
  \item \textbf{Command flags}: R (request, cleared in answers), P (proxiable), E (error answer),
        T (potentially retransmitted after a failover).
  \item \textbf{Command Code} (24 bits) identifies the command; request and answer share the code.
  \item \textbf{Application-ID} (32 bits) identifies the application (0 for the base protocol).
  \item \textbf{Hop-by-Hop Identifier}: matches answers to requests on one hop; each agent replaces it.
  \item \textbf{End-to-End Identifier}: unique per request end to end; used to detect duplicates.
\end{itemize}
Each AVP has a 32-bit code, flags V (vendor-specific: a Vendor-ID follows; 10415 for 3GPP), M (mandatory:
the receiver must understand it or reject the message) and P (end-to-end security, deprecated), a 24-bit
length that covers header and data but not padding, and data padded to a 4-byte boundary. Data types
include Integer32/64, Unsigned32/64, OctetString, UTF8String, DiameterIdentity, Address, Enumerated and
\emph{Grouped} -- an AVP whose data is itself a sequence of AVPs, which lets Diameter express nested
structures such as a subscription profile.

\begin{acnfigure}
  \acnfig{lecture-10/header}
  \caption{Diameter header (left) and AVP (right). Each row is 32 bits.}
  \label{fig:header}
\end{acnfigure}

\subsection{Commands and result codes}
Commands are named with a pair of abbreviations ending in R (request) and A (answer). The base protocol
defines peer management -- \textbf{CER/CEA} (Capabilities-Exchange, code 257), \textbf{DWR/DWA}
(Device-Watchdog, 280), \textbf{DPR/DPA} (Disconnect-Peer, 282) -- and session commands such as
\textbf{RAR/RAA} (Re-Auth, 258), \textbf{STR/STA} (Session-Termination, 275) and \textbf{ASR/ASA}
(Abort-Session, 274). Answers carry a \texttt{Result-Code} AVP whose thousands digit gives the class (\cref{tab:rc}).
\begin{table}[htbp]
\centering\small
\caption{Diameter result-code classes with examples \cite{rfc6733}.}\label{tab:rc}
\rowcolors{2}{acnmist}{white}
\begin{tabularx}{\linewidth}{@{}l l Y@{}}
\toprule
\rowcolor{white}Class & Meaning & Examples \\
\midrule
1xxx & Informational          & 1001 multi-round authentication \\
2xxx & Success                & 2001 DIAMETER\_SUCCESS \\
3xxx & Protocol errors        & 3002 UNABLE\_TO\_DELIVER, 3004 TOO\_BUSY (sent with the E flag) \\
4xxx & Transient failures     & 4001 AUTHENTICATION\_REJECTED -- may succeed later \\
5xxx & Permanent failures     & 5001 AVP\_UNSUPPORTED, 5012 UNABLE\_TO\_COMPLY \\
\bottomrule
\end{tabularx}
\end{table}
Applications define their own codes in an \texttt{Experimental-Result} AVP together with a Vendor-ID; for
example 3GPP's S6a uses 5001 for \emph{user unknown} and 5420 for \emph{unknown EPS subscription}
\cite{ts29272}. The same number therefore means different things in Result-Code and in a 3GPP
Experimental-Result -- a frequent source of confusion when reading traces.

\subsection{Diameter applications in LTE}
\begin{table}[htbp]
\centering\small
\caption{Main 3GPP Diameter applications in the EPC.}\label{tab:apps}
\rowcolors{2}{acnmist}{white}
\begin{tabularx}{\linewidth}{@{}l l l Y@{}}
\toprule
\rowcolor{white}Interface & Between & Application-ID & Main commands \\
\midrule
S6a/S6d & MME/SGSN -- HSS & 16777251 & AIR/AIA (318), ULR/ULA (316), CLR (317), IDR (319), DSR (320), PUR (321), NOR (323) \\
S13     & MME -- EIR      & 16777252 & ECR/ECA (324): ME identity check \\
Gx      & PCEF -- PCRF    & 16777238 & CCR/CCA (272), RAR/RAA (258) \\
Rx      & AF -- PCRF      & 16777236 & AAR/AAA (265), STR/STA, RAR/RAA, ASR/ASA \\
Gy      & PCEF -- OCS     & 4 (credit control) & CCR/CCA (272): INITIAL, UPDATE, TERMINATION \\
Cx/Dx   & CSCF -- HSS     & 16777216 & UAR, SAR, MAR, LIR (IMS registration) \\
\bottomrule
\end{tabularx}
\end{table}

\paragraph{S6a} (\cref{fig:s6a}) \cite{ts29272}: the MME obtains vectors with AIR/AIA and registers with
ULR/ULA, which returns the subscription data (allowed APNs with QoS profiles, AMBR, access restrictions).
The HSS uses CLR to remove the subscriber from a previous MME and IDR/DSR to change the stored profile; the
MME uses PUR when it purges an inactive subscriber.
\paragraph{Gx} \cite{ts29212}: at PDN connection set-up the P-GW sends CCR-Initial and receives PCC rules in
CCA-Initial; the PCRF later pushes changes with RAR; CCR-Update reports events (e.g.\ usage thresholds);
CCR-Termination ends the session.
\paragraph{Rx} \cite{ts29214}: the P-CSCF describes a call's media in AAR; the PCRF derives QoS and installs
rules over Gx (\cref{fig:volte}); STR ends the media session.
\paragraph{Gy} credit control \cite{rfc8506}: the PCEF requests a quota (time or volume) for a rating group in
CCR-Initial, reports usage and requests more with CCR-Update, and returns unused quota with
CCR-Termination; if credit is exhausted the OCS answers with a final-unit indication.

\subsection{Diameter in 5G}
The 5G core replaces Diameter by HTTP/2 service-based interfaces (Lectures~18--19): S6a $\to$ N8 (Nudm),
Gx $\to$ N7 (Npcf\_SMPolicyControl), Rx $\to$ N5 (Npcf\_PolicyAuthorization), Gy/Gz $\to$ N40 (Nchf). Reasons:
(i) web technology and tooling (REST, JSON, OpenAPI) are ubiquitous; (ii) network functions register and
discover each other dynamically (NRF) instead of relying on static peer tables; (iii) stateless,
cloud-native implementations scale more easily with HTTP load balancers. During the 4G/5G coexistence,
interworking functions translate between the two worlds, and the HSS and UDM are often a single
combined product.

% ---------------------------------------------------------------------
\section{Mathematical models}

\subsection{AVP and message length}\label{sec:len}
\paragraph{Rule} For an AVP with $\ell$ bytes of data, header length $h = 8$ (or 12 with Vendor-ID):
\begin{equation}\label{eq:avplen}
  L_{\mathrm{AVP}} = h + \ell, \qquad
  L_{\mathrm{on\,wire}} = 4\left\lceil \frac{h+\ell}{4}\right\rceil,
\end{equation}
and the message length is $20 + \sum L_{\mathrm{on\,wire}}$ over all top-level AVPs.
\paragraph{Example} Origin-Host \texttt{mme01.epc.mnc010.mcc432.3gppnetwork.org} has 39 characters:
$L_{\mathrm{AVP}} = 8+39 = 47$, padded to 48 bytes on the wire. Origin-Realm
\texttt{epc.mnc010.mcc432.3gppnetwork.org} (33 characters): 41, padded to 44. Identity AVPs alone cost
about 100 bytes per message, which is why Diameter messages are typically several hundred bytes even for
simple requests.

\subsection{Failure detection with the watchdog}\label{sec:watchdog}
A peer sends DWR when no traffic has been received for $T_\mathrm{w}$ (recommended 30\,s, with random jitter
of $\pm2$\,s to avoid synchronisation). If no answer arrives within a further $T_\mathrm{w}$, the connection
is considered suspect and traffic is failed over to an alternative peer. The worst-case detection time on
an idle connection is therefore about
\begin{equation}\label{eq:tdet}
  T_{\mathrm{det}} \approx 2\,T_\mathrm{w},
\end{equation}
i.e.\ about a minute with defaults; on a busy connection, unanswered application requests reveal the failure
much sooner. Operators tune $T_\mathrm{w}$ down (e.g.\ 6\,s) on critical links, trading watchdog traffic for
faster failover (\cref{exr:watchdog}).

\subsection{HSS transaction load}\label{sec:hssload}
\paragraph{Model} $N$ subscribers each attach $a$ times per day and perform $u$ MME-changing TAUs per day,
each needing one AIR and one ULR (AIR only if the MME lacks vectors); other procedures are neglected. The
mean S6a transaction rate is
\begin{equation}\label{eq:hss}
  \lambda_{\mathrm{S6a}} \approx \frac{N\,(a+u)\,\bar n}{86\,400}\ \text{transactions/s},
\end{equation}
with $\bar n$ the mean number of transactions per procedure.
\paragraph{Example} $N = 10^7$, $a+u = 3$, $\bar n = 2$: $\lambda \approx 694$\,transactions/s on average. The
busy hour is several times higher, and a restart of an MME pool can create bursts far above the mean
(Lecture~9), so HSS front-ends are dimensioned for peaks, not averages.

\subsection{Credit-control quota sizing}\label{sec:quota}
\paragraph{Problem} How large should a Gy volume quota be?
\paragraph{Model} A user consumes data at rate $r$ (bytes/s); each grant gives $Q$ bytes; the PCEF reports
and requests a new quota when the current one is used.
\paragraph{Results} The number of CCR-Updates per user per unit time and the maximum overspend (consumption
the operator cannot charge if the user's credit ran out just after the grant) are
\begin{equation}\label{eq:quota}
  \lambda_{\mathrm{CCR}} = \frac{r}{Q}, \qquad O_{\max} \approx Q .
\end{equation}
\paragraph{Example and interpretation} $r = \qty{2}{\mega\bit\per\second} = 250$\,kB/s, $Q = 10$\,MB: one update
every 40\,s (90 per hour) and up to 10\,MB of revenue at risk per user. Doubling $Q$ halves the signalling
but doubles the risk -- operators therefore give small quotas to low-balance users and large ones to
high-balance users. A validity time limits how long a quota may be held.

% ---------------------------------------------------------------------
\section{Network architecture and diagrams}
In a typical operator network, every MME, P-GW, PCRF, HSS and OCS connects to a redundant pair of DRAs;
the DRAs connect to each other and to a pair of DEAs facing the roaming exchange. Routing is by
realm and application, with load balancing and session binding (all messages of a Gx/Rx session must
reach the same PCRF instance, which DRAs ensure by binding the UE's IP address to the PCRF). Diameter is
purely control plane: none of these nodes ever forwards user packets.

\begin{acnfigure}
  \acnfig{lecture-10/peer-lifecycle}
  \caption{Diameter peer lifecycle: transport set-up, capabilities exchange, watchdog, disconnection.}
  \label{fig:peer}
\end{acnfigure}

% ---------------------------------------------------------------------
\section{Step-by-step procedures}

\begin{procedurebox}{S6a during attach}\label{proc:s6a}
\begin{steps}
  \item The MME needs vectors for an unknown IMSI. It builds AIR (code 318, Application-ID 16777251,
        R flag set) with Session-Id, Auth-Session-State, Origin-Host/Realm, Destination-Realm derived from the
        IMSI, User-Name = IMSI, Visited-PLMN-Id and Requested-EUTRAN-Authentication-Info.
  \item The DRA routes on Destination-Realm and Application-ID to an HSS front-end; the HSS returns AIA with
        Result-Code 2001 and Authentication-Info (grouped E-UTRAN-Vector AVPs).
  \item After authentication, the MME sends ULR (316) with RAT-Type and ULR-Flags; the HSS stores the MME
        identity, sends CLR to the previous MME if there was one, and returns ULA with Subscription-Data.
  \item Errors appear as Experimental-Result, e.g.\ 5001 (user unknown): the MME rejects the attach with an
        appropriate NAS cause.
\end{steps}
\end{procedurebox}

\begin{acnfigure}
  \resizebox{\linewidth}{!}{\acnfig{lecture-10/s6a-attach}}
  \caption{S6a exchanges during attach via a Diameter Routing Agent (\cref{proc:s6a}).}
  \label{fig:s6a}
\end{acnfigure}

\begin{procedurebox}{Policy for a VoLTE call (Rx + Gx)}\label{proc:volte}
\begin{steps}
  \item The P-CSCF, having seen the SDP offer/answer of the call, sends Rx AAR to the PCRF with
        Media-Component-Description AVPs (codec, bandwidth, flow descriptions).
  \item The PCRF authorises the media, derives QCI~1 and GBR, answers AAA and sends Gx RAR with
        Charging-Rule-Install to the P-GW serving the UE's IMS PDN connection.
  \item The P-GW creates a dedicated bearer (Lecture~9) and answers RAA.
  \item At call end, the P-CSCF sends STR; the PCRF removes the rules (RAR with Charging-Rule-Remove), and
        the bearer is deleted.
\end{steps}
\end{procedurebox}

\begin{acnfigure}
  \resizebox{\linewidth}{!}{\acnfig{lecture-10/volte-policy}}
  \caption{Rx and Gx dialogues for a VoLTE call (\cref{proc:volte}).}
  \label{fig:volte}
\end{acnfigure}

% ---------------------------------------------------------------------
\section{Worked examples}

\begin{examplebox}{Decoding a header}\label{ex:decode}
\textbf{Problem.} The first 20 bytes of a message are (hex)
\texttt{01 00 01 2C 80 00 01 3C 01 00 00 23 \ldots}. Decode version, length, flags, command and application.

\textbf{Solution.} Version 1; length $\mathtt{0x00012C} = 300$ bytes; flags $\mathtt{0x80}$ = R only (a request,
not proxiable); command $\mathtt{0x00013C} = 316$ = Update-Location; Application-ID $\mathtt{0x01000023} =
16777251$ = S6a. So: a ULR from an MME. (Real S6a requests usually also set P.)
\end{examplebox}

\begin{examplebox}{Message size}\label{ex:size}
\textbf{Problem.} An AIR contains: Session-Id (60 bytes data), Auth-Session-State (Unsigned32), Origin-Host and
Origin-Realm as in \cref{sec:len}, Destination-Realm (33 bytes), User-Name (15-digit IMSI), Visited-PLMN-Id
(3-byte OctetString, vendor-specific) and Requested-EUTRAN-Authentication-Info (grouped, vendor-specific,
containing one Number-Of-Requested-Vectors Unsigned32 AVP, vendor-specific). Compute its length.

\textbf{Solution.} Session-Id $4\lceil68/4\rceil = 68$; Auth-Session-State $8+4 = 12$; Origin-Host 48;
Origin-Realm 44; Destination-Realm 44; User-Name $4\lceil23/4\rceil = 24$; Visited-PLMN-Id $4\lceil15/4\rceil
= 16$; inner AVP $12+4 = 16$, so the grouped AVP is $12+16 = 28$. Total $20 + 68+12+48+44+44+24+16+28 = 304$
bytes.
\end{examplebox}

\begin{examplebox}{Quota policy}\label{ex:quota}
\textbf{Problem.} An OCS must handle at most 2000 CCR-Updates per second for $10^6$ active data users with a
mean rate of 50\,kB/s. What is the smallest uniform quota?

\textbf{Solution.} $\lambda_{\mathrm{total}} = N r/Q \le 2000 \Rightarrow Q \ge 10^6\times5\times10^4/2000 =
25$\,MB. The operator would use 25\,MB (or more) and limit risk with validity timers and smaller quotas for
low-balance users.
\end{examplebox}

% ---------------------------------------------------------------------
\section{Technology comparison}
\begin{table}[htbp]
\centering\small
\caption{Core signalling protocols across generations.}\label{tab:sig}
\rowcolors{2}{acnmist}{white}
\begin{tabularx}{\linewidth}{@{}l Y Y Y@{}}
\toprule
\rowcolor{white} & SS7 MAP (2G/3G) & Diameter (4G, IMS) & HTTP/2 SBI (5G) \\
\midrule
Transport    & MTP/SCCP (or SIGTRAN over IP) & TCP/SCTP, port 3868 & TCP + TLS, HTTP/2 \\
Encoding     & ASN.1 BER                     & Binary AVPs (TLV)   & JSON (OpenAPI schemas) \\
Addressing   & Global titles, point codes    & Realms, hosts       & URIs, NRF discovery \\
Routing      & STPs                          & DRAs                & Direct or via SCP \\
Roaming edge & International STP             & DEA                 & SEPP \\
Interaction  & TCAP dialogues                & Request/answer sessions & RESTful resources, subscribe/notify \\
\bottomrule
\end{tabularx}
\end{table}

\paragraph{Diameter is control plane only} Diameter carries decisions about sessions -- who is allowed, with
which QoS, how much credit -- but no session data. The packets of the subscriber flow through GTP-U
tunnels (Lecture~13) that Diameter never touches; a failed PCRF blocks new policy decisions but not the
packets of existing sessions.

% ---------------------------------------------------------------------
\section{Key takeaways}
\begin{takeawaybox}
\begin{itemize}
  \item Diameter = reliable (TCP/SCTP), extensible AAA: base protocol + applications identified by
        Application-ID.
  \item 20-byte header (version, length, R/P/E/T flags, command code, application ID, hop-by-hop and
        end-to-end IDs) + AVPs padded to 32 bits.
  \item Peers: CER/CEA, DWR/DWA, DPR/DPA; routing by Destination-Realm via relays, proxies and DRAs/DEAs.
  \item LTE applications: S6a (AIR, ULR, CLR, IDR), Gx (CCR/CCA, RAR), Rx (AAR), Gy (credit control).
  \item Result-Code classes 1xxx--5xxx; 3GPP-specific errors in Experimental-Result.
  \item Quota sizing trades signalling $r/Q$ against revenue risk $\approx Q$.
  \item 5G replaces Diameter by HTTP/2 SBIs; interworking during 4G/5G coexistence.
\end{itemize}
\end{takeawaybox}

% ---------------------------------------------------------------------
\section{Exercises}

\begin{exercisebox}{1}{Server-initiated messages}\label{exr:srv}
Give two EPC situations in which the server must initiate a Diameter request, and explain why RADIUS
could not support them natively.
\end{exercisebox}

\begin{exercisebox}{1}{Two identifiers}\label{exr:ids}
Explain why Diameter needs both a hop-by-hop and an end-to-end identifier. What happens to each when a relay
forwards a request?
\end{exercisebox}

\begin{exercisebox}{2}{AVP length}\label{exr:avp}
Compute the on-wire length of: (a) a Result-Code AVP; (b) a vendor-specific UTF8String AVP containing the APN
\texttt{internet}; (c) a grouped AVP containing (a) and (b).
\end{exercisebox}

\begin{exercisebox}{2}{Watchdog tuning}\label{exr:watchdog}
A DRA has 400 idle peer connections. Compute the watchdog message rate for $T_\mathrm{w} = 30$\,s and $6$\,s and
the corresponding worst-case failure-detection times.
\end{exercisebox}

\begin{exercisebox}{3}{Agent topology}\label{exr:topo}
An operator has 20 MMEs, 10 HSS front-ends, 30 P-GWs, 8 PCRFs and 4 OCS nodes. Compare the number of Diameter
peer connections for a full mesh of the interfaces S6a, Gx and Gy and for a design with a redundant pair of DRAs
(every node connects to both DRAs). Discuss the reliability implications.
\end{exercisebox}

\begin{exercisebox}{3}{Quota with risk budget}\label{exr:risk}
A user with balance $B$ is granted quotas $Q = \min(Q_{\max}, \alpha B)$. Show that the expected number of
CCR-Updates to consume the whole balance at constant rate is about $\log(B/Q_{\min})/\log(1/(1-\alpha))$ when
$Q_{\max}$ is never reached and quotas stop at $Q_{\min}$, and evaluate it for $\alpha = 0.2$,
$B = 1$\,GB, $Q_{\min} = 1$\,MB.
\end{exercisebox}

\subsection*{Solutions}
\addcontentsline{toc}{subsection}{Solutions}

\begin{solutionbox}{\ref{exr:srv}}
(1) HSS $\to$ MME Cancel-Location (CLR) when the subscriber registers in another MME or is withdrawn; (2) PCRF
$\to$ P-GW Re-Auth (RAR) when a VoLTE call needs a dedicated bearer or a usage limit is reached; also HSS
Insert-Subscriber-Data when a profile changes. RADIUS clients send requests and servers only answer; there is
no standard way for a server to start a transaction (only later extensions such as Dynamic Authorization), and
NAT/firewalls in front of UDP clients make it unreliable.
\end{solutionbox}

\begin{solutionbox}{\ref{exr:ids}}
The hop-by-hop identifier lets a node match an answer to the request it sent on a given connection; each relay
replaces it with its own value and restores the original when returning the answer. The end-to-end identifier
is set by the originator and never changed; together with Origin-Host it lets the server detect duplicates
(e.g.\ a request retransmitted with the T flag after a failover) for a period of time.
\end{solutionbox}

\begin{solutionbox}{\ref{exr:avp}}
(a) $8+4 = 12$ bytes. (b) $12 + 8 = 20$ bytes (already aligned). (c) $8$ (or 12 if vendor-specific) $+12+20 =
40$ bytes (44 with a Vendor-ID).
\end{solutionbox}

\begin{solutionbox}{\ref{exr:watchdog}}
Each idle connection produces one DWR and one DWA per $T_\mathrm{w}$ from each side that initiates; counting one DWR
and one DWA per connection per $T_\mathrm{w}$: $400\times2/30 = 26.7$ messages/s for 30\,s, $133$ messages/s for
6\,s -- negligible either way. Detection time $\approx 2T_\mathrm{w}$: 60\,s versus 12\,s. The shorter timer is clearly
worthwhile on critical links; the risk is false failover on congested links.
\end{solutionbox}

\begin{solutionbox}{\ref{exr:topo}}
Full mesh: S6a $20\times10 = 200$, Gx $30\times8 = 240$, Gy $30\times4 = 120$: 560 connections, each to be
configured and monitored. With two DRAs: $(20+10+30+8+4)\times2 = 144$ connections (plus one between the DRAs).
DRAs centralise routing, load balancing and session binding and simplify adding nodes, but they become critical
single points: they must be deployed in geo-redundant pairs with enough capacity to carry all traffic if one
fails.
\end{solutionbox}

\begin{solutionbox}{\ref{exr:risk}}
Each grant uses a fraction $\alpha$ of the remaining balance, so after $k$ grants the balance is $B(1-\alpha)^k$;
quotas fall below $Q_{\min}$ when $\alpha B(1-\alpha)^k = Q_{\min}$, i.e.\
$k \approx \log(\alpha B/Q_{\min})/\log(1/(1-\alpha))$, essentially the stated expression. For $\alpha=0.2$,
$B = 1000$\,MB, $Q_{\min} = 1$\,MB: $k \approx \ln(200)/\ln(1.25) = 5.30/0.223 = 23.7$ grants, after which a few
minimum-size grants finish the balance. Geometric quotas keep both signalling (logarithmic in $B$) and risk
(at most $\alpha$ of the balance) small.
\end{solutionbox}

% ---------------------------------------------------------------------
\section{Further reading}
The base protocol is RFC~6733 \cite{rfc6733}, credit control RFC~8506 \cite{rfc8506}; RADIUS is RFC~2865
\cite{rfc2865}. The 3GPP applications are specified in \cite{ts29272} (S6a/S6d/S13), \cite{ts29212} (Gx) and
\cite{ts29214} (Rx); realm naming follows \cite{ts23003}. Olsson et al.\ \cite{olsson2013epc} discuss Diameter
deployment in the EPC.

\printbibliography[heading=subbibliography,title={\ACNstrRefs}]
\end{document}
